\section{Discussion}
\label{sec:discussion}

\para{What do the results imply?} The central empirical answer is negative to the strongest form of the null hypothesis: hidden states in different MLP layers are not totally different. The clearest evidence is local. Adjacent layers beat non-adjacent layers in every within-model comparison, often by large margins under \cka, \svcca, and \pwcca. At the same time, deeper layers improve probe accuracy, which means similarity is not merely redundancy. The more plausible interpretation is that depth progressively refines a shared representational space.

\para{Why are cross-seed results weaker?} Cross-seed alignment asks more than continuity. It asks whether independently trained models reuse the same depth-specific organization. The mixed results suggest that a layer can play a similar functional role without preserving one stable geometric template across seeds. This is especially clear on \cifar, where \pwcca and \svcca become flat or negative despite positive within-model locality. For practical interpretability, this means layer-to-layer comparison is safer inside one model than across separately trained replicas.

\para{Limitations.} This is a controlled study of plain feedforward MLPs on flattened image inputs. The findings should not be transferred directly to architectures with convolution, residual pathways, or attention. The runtime environment also forced CPU execution and moderate dataset subsampling despite four visible GPUs, which limits scale and may suppress stronger convergence on \cifar. Finally, every similarity score compresses rich geometry into a scalar summary, and \svcca/\pwcca can be numerically sensitive.

\para{Broader implications.} Even with those limitations, the study provides a useful empirical baseline. It supports the idea that layer-wise analysis in simple feedforward networks is meaningful, and it shows why multi-metric reporting is necessary. The most direct next extension is functional: layer stitching could test whether the observed geometric continuity translates into interchangeable computation. More broadly, the same benchmark design can be used to study when progressive similarity breaks down as width, depth, or architecture changes.
